\section{Kuasarlar ve Mikrokuasarlar}

Kuasarlar ve mikrokuasarlar, astrofizikte \emph{akresyonla beslenen kompakt cisimler}in en uç örnekleridir. Her iki sınıf da yüksek enerjili ışınım, relativistik jetler ve güçlü yerçekimi alanları ile karakterize edilir. Aralarındaki temel fark, merkezi kara deliğin kütle ölçeği ve sistemin kozmolojik bağlamıdır.

\subsection{Kuasar Nedir?}

Kuasar (Quasi-Stellar Radio Source), bir galaksinin merkezinde bulunan \emph{süper kütleli kara delik} (Supermassive Black Hole, SMBH) tarafından beslenen, son derece parlak aktif galaktik çekirdeklerdir (AGN).

Tipik kara delik kütlesi:
\begin{equation}
M_{\mathrm{BH}} \sim 10^6\text{--}10^{10}\,M_\odot
\end{equation}

Kuasarların toplam elektromanyetik parlaklığı:
\begin{equation}
L \sim 10^{45}\text{--}10^{48}\ \mathrm{erg\,s^{-1}}
\end{equation}
olup, tüm evrendeki en parlak sürekli kaynaklar arasındadır.

\subsubsection{Enerji Kaynağı: Akresyon}

Kuasarların enerji kaynağı, kara deliğe düşen maddenin kütleçekim potansiyel enerjisidir. Akresyon diski için yayılan toplam parlaklık:
\begin{equation}
L = \eta \dot{M} c^2
\end{equation}
burada $\eta \sim 0.06$--$0.42$ akresyon verimidir (kara deliğin dönmesine bağlıdır).

\subsubsection{Eddington Sınırı}

Kuasar parlaklığı çoğunlukla Eddington parlaklığı ile karşılaştırılır:
\begin{equation}
L_{\mathrm{Edd}} = \frac{4\pi G M m_p c}{\sigma_T}
\end{equation}

Bir kuasar genellikle:
\begin{equation}
L \sim (0.1\text{--}1)\,L_{\mathrm{Edd}}
\end{equation}
rejiminde çalışır.

\subsubsection{Relativistik Jetler}

Birçok kuasar, ışık hızına yakın hızlarda fırlatılan çift kutuplu jetler üretir. Jet Lorentz faktörü:
\begin{equation}
\Gamma = \frac{1}{\sqrt{1 - v^2/c^2}} \sim 5\text{--}20
\end{equation}

Jetler, senkrotron ve ters Compton süreçleriyle radyo–gamma bantlarında ışınım yapar.

\subsection{Mikrokuasar Nedir?}

Mikrokuasarlar, \emph{Galaksi içinde} bulunan ve kuasarların küçük ölçekli analogları olan sistemlerdir. Bir mikrokuasar, genellikle:
\begin{itemize}
    \item Yıldız kütleli bir kara delik veya nötron yıldızı
    \item Bir eş yıldız
    \item Akresyon diski
    \item Relativistik jetler
\end{itemize}
içeren bir X-ışını ikilisidir.

Tipik merkezi cisim kütlesi:
\begin{equation}
M_{\mathrm{compact}} \sim 5\text{--}20\,M_\odot
\end{equation}

\subsubsection{Neden ``Mikro''? Ölçekleme Yasaları}

Akresyon fiziği, boyutsuz değişkenlerle yazıldığında kütleyle ölçeklenir. Zaman ölçeği:
\begin{equation}
t \propto \frac{GM}{c^3}
\end{equation}

Bu nedenle kuasarlarda binlerce yıl süren süreçler, mikrokuasarlarda:
\begin{equation}
t_{\mu Q} \sim \left(\frac{M_{\mu Q}}{M_{\mathrm{QSO}}}\right) t_{\mathrm{QSO}}
\end{equation}
ölçeğiyle saniye–gün aralığında gözlenebilir.

\subsubsection{X-Işını ve Jet Yayınımı}

Mikrokuasarlarda:
\begin{itemize}
    \item X-ışını: İç akresyon diski ve korona
    \item Radyo/IR: Jet senkrotron yayınımı
\end{itemize}
baskındır.

Jet hızı:
\begin{equation}
v_{\mathrm{jet}} \sim (0.8\text{--}0.99)c
\end{equation}

\subsection{Kuasar ve Mikrokuasar Karşılaştırması}

\begin{center}
\begin{tabular}{lcc}
\hline
 & Kuasar & Mikrokuasar \\
\hline
Merkezi cisim & SMBH & Yıldız kütleli BH/NS \\
Kütle & $10^6$--$10^{10} M_\odot$ & $5$--$20 M_\odot$ \\
Konum & Galaksi merkezi & Galaksi içi \\
Zaman ölçeği & $10^3$--$10^6$ yıl & saniye--gün \\
Jetler & Var & Var \\
\hline
\end{tabular}
\end{center}

\subsection{Astrofiziksel Önemi}

Kuasarlar:
\begin{itemize}
    \item Erken evrenin aydınlatılması
    \item Galaksi–kara delik birlikte evrimi
    \item Kozmik yeniden iyonlaşma
\end{itemize}

Mikrokuasarlar:
\begin{itemize}
    \item Akresyon ve jet fiziğinin zamana bağlı incelenmesi
    \item Görelilik etkilerinin doğrudan test edilmesi
    \item Kuasarların laboratuvar ölçekli analogları
\end{itemize}

Bu iki sınıf, evrensel akresyon fiziğinin farklı kütle ölçeklerindeki tezahürleridir.
